% Experiment Designs poster — NEMI 2026, v3
% Six criteria grounded in the mechanistic validity paper.
% All designs properly attributed, claims calibrated.
% Wording reviewed against hostile-review feedback.
%
%   cd poster && lualatex designs_poster_v3.tex

\documentclass[25pt, a0paper, portrait]{tikzposter}

\usepackage[T1]{fontenc}
\usepackage{amsmath,amssymb}
\usepackage{array}
\usepackage{booktabs}
\usepackage{enumitem}
\usepackage{microtype}

\tikzposterlatexaffectionproofoff

% ── Color palette ───────────────────────────────────────
\definecolor{darktext}{HTML}{2C3E50}
\definecolor{blockbg}{HTML}{FFFFFF}
\definecolor{hookred}{HTML}{C0392B}
\definecolor{designblue}{HTML}{2980B9}
\definecolor{constructgreen}{HTML}{27AE60}
\definecolor{externalpurple}{HTML}{8E44AD}

% ── Theme ───────────────────────────────────────────────
\usetheme{Default}

\definecolorstyle{DesignColors}{
  \definecolor{colorOne}{HTML}{2C3E50}
  \definecolor{colorTwo}{HTML}{3498DB}
  \definecolor{colorThree}{HTML}{ECF0F1}
}{
  \colorlet{backgroundcolor}{colorThree}
  \colorlet{framecolor}{colorOne}
  \colorlet{titlefgcolor}{white}
  \colorlet{titlebgcolor}{colorOne}
  \colorlet{blocktitlebgcolor}{designblue}
  \colorlet{blocktitlefgcolor}{white}
  \colorlet{blockbodybgcolor}{blockbg}
  \colorlet{blockbodyfgcolor}{darktext}
  \colorlet{notefgcolor}{darktext}
  \colorlet{notebgcolor}{designblue!12}
  \colorlet{notefrcolor}{designblue}
}
\usecolorstyle{DesignColors}

\defineblockstyle{DesignBlock}{
  titlewidthscale=1.0, bodywidthscale=1.0, titleleft,
  titleoffsetx=0pt, titleoffsety=0pt,
  bodyoffsetx=0pt, bodyoffsety=0pt, bodyverticalshift=0pt,
  roundedcorners=12, linewidth=2pt, titleinnersep=10pt, bodyinnersep=14pt
}{
  \draw[rounded corners=\blockroundedcorners, color=blocktitlebgcolor,
        fill=blockbodybgcolor, line width=\blocklinewidth]
    (blockbody.south west) rectangle (blocktitle.north east);
  \ifBlockHasTitle
    \draw[rounded corners=\blockroundedcorners,
          fill=blocktitlebgcolor, color=blocktitlebgcolor,
          line width=\blocklinewidth]
      (blocktitle.south west) rectangle (blocktitle.north east);
  \fi
}
\useblockstyle{DesignBlock}

\definetitlestyle{DesignTitle}{
  width=\paperwidth, roundedcorners=0, linewidth=0pt, innersep=20pt,
  titletotopverticalspace=15mm, titletoblockverticalspace=25mm
}{
  \draw[fill=titlebgcolor, color=titlebgcolor]
    (\titleposleft, \titleposbottom)
    rectangle (\titleposright, \titlepostop);
}
\usetitlestyle{DesignTitle}

\title{%
  \parbox{0.92\linewidth}{\centering\bfseries
    What Does Your Ablation Rule Out?}}
\author{Elliot Tower\quad\texttt{elliot@elliottower.ai}}
\institute{}

\begin{document}

\maketitle[width=\paperwidth]

% ════════════════════════════════════════════════════════
% ROW 1 — THE PITCH + COMPLEMENTATION (C6)
% ════════════════════════════════════════════════════════
\begin{columns}
\column{0.45}

{
\colorlet{blocktitlebgcolor}{hookred}
\block{Different Designs Answer Different Causal Questions}{
  Ablating a circuit and measuring degradation establishes causal
  involvement. Stronger conclusions---specificity, independence,
  reversibility, training-time coupling---require additional designs,
  each of which rules out a different rival explanation.

  \vspace{5mm}
  The ORBITA trial (2017) illustrates the general principle: compare
  an intervention against a control that preserves its incidental
  features while removing the hypothesized active ingredient. In MI,
  this motivates matched perturbation controls rather than comparison
  with an intact model alone.

  \vspace{5mm}
  These designs are not unknown in MI. The claim is that they are not
  yet routine, not consistently interpreted as distinct inferential
  tests, and not standardized in reporting.

  \vspace{5mm}
  \begin{center}
  \renewcommand{\arraystretch}{1.3}
  {\small
  \begin{tabular}{@{}llll@{}}
    \toprule
    \textbf{ID} & \textbf{Design} & \textbf{Rules out} & \textbf{Cost} \\
    \midrule
    C6  & Complementation     & Mislabeled units   & 2 passes \\
    I6  & Double dissociation & Generic disruption  & 1 cond. \\
    I9  & Factorial ablation  & Hidden interaction  & 1 cond. \\
    I10 & Rescue              & Collateral damage   & 1 pass \\
    I11/12 & Onset/offset     & Spurious coexistence & Varies \\
    E5  & Dose--response      & Threshold artifacts & Low \\
    \bottomrule
  \end{tabular}}
  \end{center}
}
}

\column{0.55}

{
\colorlet{blocktitlebgcolor}{constructgreen}
\block{C6: Complementation --- Same Unit or Different?}{
  In genetics, two mutations that both break a function might be in
  the same gene or in different genes. Benzer's cis-trans test (1955)
  distinguishes them: place one mutation on each chromosome (the
  \emph{trans} configuration). If function recovers, the mutations
  are in different genes---they \emph{complement}.

  \vspace{5mm}
  \textbf{MI translation.} S-inhibition heads and negative name
  movers both suppress the repeated name in IOI. Run two forward
  passes, one with each group intact. Construct a hybrid activation
  by taking the S-inhibition contribution from pass~1 and the
  negative name-mover contribution from pass~2 (trans configuration).
  If the model recovers, they serve independent functional roles. If
  the hybrid fails, they are part of the same functional unit.

  \vspace{5mm}
  \begin{center}
  \renewcommand{\arraystretch}{1.3}
  {\large
  \begin{tabular}{@{}lll@{}}
    \toprule
    \textbf{Configuration} & \textbf{Result} & \textbf{Interpretation} \\
    \midrule
    Trans (hybrid pass) & Recovers & Different units \\
    Trans (hybrid pass) & Fails    & Same unit \\
    Cis (single pass)   & Works    & Control \\
    \bottomrule
  \end{tabular}}
  \end{center}

  \vspace{5mm}
  \textbf{Rules out:} treating two labels as two mechanisms without
  testing whether they name the same functional unit.

  \vspace{4mm}
  Across 16 audited claims, this test has been applied
  \textbf{zero times}. Circuit papers identify component classes
  without testing whether those classes are functionally distinct.

  \vspace{4mm}
  \textbf{Cost:} two forward passes + hybrid assembly.

  \vspace{3mm}
  {\footnotesize Genetics. Benzer's cis-trans test (1955).}
}
}

\end{columns}

% ════════════════════════════════════════════════════════
% ROW 2 — THREE INTERNAL/EXTERNAL CRITERIA
% ════════════════════════════════════════════════════════
\begin{columns}
\column{0.33}

\block{I6: Double Dissociation}{
  A selective effect---one intervention disrupts one behavior---shows
  involvement. A \textbf{double dissociation} shows specificity: two
  interventions with crossed effects.

  \vspace{4mm}
  \begin{center}
  \renewcommand{\arraystretch}{1.3}
  \begin{tabular}{@{}lcc@{}}
    \toprule
    & \textbf{Behav.\ X} & \textbf{Behav.\ Y} \\
    \midrule
    Ablate mech.\ A & Large & Small \\
    Ablate mech.\ B & Small & Large \\
    \bottomrule
  \end{tabular}
  \end{center}

  \vspace{4mm}
  A selective effect in one row is a \textbf{single dissociation}.
  Many MI papers achieve this---showing that an intervention
  selectively affects one condition. The crossed pattern across
  \emph{both} rows is a double dissociation. Neither effect alone
  can be explained by generalized disruption.

  \vspace{4mm}
  \textbf{Rules out:} involvement mistaken for specificity.

  \vspace{4mm}
  Across 16 audited claims, exactly \textbf{one} meets this
  criterion (Feucht et~al.\ 2025, three years after the original
  paper).

  \vspace{4mm}
  \textbf{Cost:} one extra condition.

  \vspace{3mm}
  {\footnotesize Neuroscience. Shallice (1988).}
}

\column{0.34}

\block{I9: Factorial Ablation / Epistasis}{
  Do circuit components contribute independently? Ablate A alone,
  B alone, both together. Compare the joint deficit with the sum of
  single deficits.

  \vspace{4mm}
  \begin{center}
  \renewcommand{\arraystretch}{1.3}
  \begin{tabular}{@{}ll@{}}
    \toprule
    \textbf{Pattern} & \textbf{Interpretation} \\
    \midrule
    Joint $=$ sum   & Independent \\
    Joint $<$ sum   & Redundancy or shared bottleneck \\
    Joint $>$ sum   & Synergy or mutual dependence \\
    \bottomrule
  \end{tabular}
  \end{center}

  \vspace{4mm}
  Subadditivity can indicate redundancy \emph{or} a shared
  downstream bottleneck; superadditivity can indicate synergy
  \emph{or} a ceiling effect. Neither pattern alone establishes
  mechanistic identity---the additive null must be assessed on an
  appropriate scale for bounded metrics.

  \vspace{4mm}
  \textbf{Rules out:} treating two component classes as independent
  contributors without measuring their interaction.

  \vspace{4mm}
  \textbf{Cost:} one extra ablation condition.

  \vspace{3mm}
  {\footnotesize Genetics. Fisher's epistasis; interaction structure
  in gene networks.}
}

\column{0.33}

{
\colorlet{blocktitlebgcolor}{externalpurple}
\block{E5: Dose--Response}{
  An intervention-strength sweep reveals thresholds, saturation,
  compensation, and non-monotonicity that a single ablation endpoint
  cannot show.

  \vspace{4mm}
  \textbf{MI translation.} Sweep ablation strength from 0\% to
  100\%. The shape of the resulting curve carries more information
  than either endpoint.

  \vspace{4mm}
  \textbf{Caveat.} Interpretation depends on whether intervention
  strength is a meaningful dose. Linear interpolation in activation
  space need not represent proportional mechanistic inhibition.
  Genuine mechanisms can have thresholds; nonspecific disruption
  can produce smooth monotonic decline.

  \vspace{4mm}
  \textbf{Rules out:} threshold artifacts and saturation effects
  invisible at a single intervention strength.

  \vspace{4mm}
  \textbf{Cost:} low marginal cost, but interpretation requires
  validating the intervention scale.

  \vspace{3mm}
  {\footnotesize Pharmacology. Dose-response curves, Hill (1910).}
}
}

\end{columns}

% ════════════════════════════════════════════════════════
% ROW 3 — RESCUE, ONSET/OFFSET, TAKEAWAY
% ════════════════════════════════════════════════════════
\begin{columns}
\column{0.33}

\block{I10: Rescue / Reversibility}{
  After producing a deficit by ablation, restore the hypothesized
  mediator using an independently specified intervention.

  \vspace{4mm}
  \textbf{MI translation.} Re-inject the clean activation after
  ablation and measure recovery. Activation patching already
  contains this logic: patching in a clean activation \emph{is} a
  rescue experiment. Almost no published evaluation reports the
  recovery rate alongside the degradation rate.

  \vspace{4mm}
  Successful recovery strengthens causal attribution. Failed
  recovery is \textbf{ambiguous}---the restoration may be at the
  wrong site, the intervention may create off-manifold states, or
  the component may require coordinated restoration with its inputs.

  \vspace{4mm}
  \textbf{Rules out} (when recovery succeeds): collateral damage
  from the original ablation.

  \vspace{4mm}
  \textbf{Cost:} one forward pass.

  \vspace{3mm}
  {\footnotesize Genetics. Rescue experiments in functional
  genomics.}
}

\column{0.34}

\block{I11/I12: Onset and Offset Coupling}{
  Does causal contribution emerge with the capability (onset, I11)
  and diminish when the capability is selectively removed (offset,
  I12)?

  \vspace{4mm}
  \textbf{MI translation.} Onset: track the circuit across training
  checkpoints; the mechanism should become causally active when the
  capability appears. Offset: fine-tune or unlearn the capability,
  then check whether the circuit's causal contribution diminishes.

  \vspace{4mm}
  Onset alone is reported for induction heads---the circuit forms at
  a training phase transition when in-context copying emerges. The
  partner test, offset coupling, has not been reported.

  \vspace{4mm}
  Concordant onset and offset are stronger than either alone, but
  require controls for global representation drift and mechanistic
  rerouting---a mechanism can persist dormant after its capability
  is suppressed.

  \vspace{4mm}
  \textbf{Rules out:} spurious coexistence at a single checkpoint.

  \vspace{4mm}
  \textbf{Cost:} onset = checkpoint tracking (low). Offset =
  fine-tuning or unlearning (high).

  \vspace{3mm}
  {\footnotesize Medical microbiology. Fredricks \& Relman (1996).}
}

\column{0.33}

{
\colorlet{blocktitlebgcolor}{colorOne}
\block{Takeaway}{
  {\Large\bfseries Different intervention designs answer different
  causal questions. Ablation alone does not answer all of them.}

  \vspace{5mm}
  \begin{itemize}[leftmargin=*, itemsep=3mm]
    \item \textbf{1/16} audited claims meets double
      dissociation (I6)
    \item \textbf{0/16} meets complementation (C6)
    \item \textbf{0/16} reports epistatic interaction
      structure (I9)
    \item \textbf{0/16} reports offset coupling (I12)
    \item Each design rules out a rival explanation that
      ablation leaves open
  \end{itemize}

  \vspace{4mm}
  {\small These six criteria are part of a 36-criterion validity
  framework spanning construct, measurement, internal, external, and
  interpretive validity. The framework also introduces interpretive
  validity (V1--V5)---the requirement that the description mode match
  the evidence---for which no existing validity framework has a
  counterpart.}

  \vspace{5mm}
  \begin{center}
  \texttt{elliot@elliottower.ai}

  \vspace{2mm}
  {\small Paper:}
  \texttt{doi.org/10.5281/zenodo.20478480}
  \end{center}
}
}

\end{columns}

\end{document}
